\documentclass{article}
\usepackage{amsmath,amssymb}
\usepackage[vocab={digits}]{leanify}
\newtheorem{lemma}{Lemma}
\begin{document}

\section{Notation}

We write $d_k(n)$ for the $k$th digit of $n$ in base $10$.
We write $\ell_b(n)$ for the number of digits of $n$ in base $b$.
We write $R_b(n)$ for the digit reversal of $n$ in base $b$.
We write $\rho(n)$ for the digital root of $n$.
We write $U_n$ for the $n$th repunit in base $10$.

\section{Digits}

\begin{lemma}\label{digits}
$d_0(1234)=4$, $d_2(1234)=2$ and $d_5(1234)=0$.
\end{lemma}

\begin{lemma}\label{length}
$\ell_{10}(1234)=4$ and $\ell_2(8)=4$.
\end{lemma}

\begin{lemma}\label{reversal}
$R_{10}(1234)=4321$, and $121$ is a palindrome in base $10$ and $5$ is a palindrome in base $2$.
\end{lemma}

\begin{lemma}\label{root}
$\rho(1234)=1$ and $U_3=111$.
\end{lemma}

\begin{lemma}\label{special}
$18$ is a Harshad number and $153$ is a narcissistic number.
\end{lemma}

\begin{lemma}\label{harshad}
Every Harshad number is a positive integer.
\end{lemma}

\end{document}
